\section{Methodology}
\label{sec:method}

In this section, we present the formal mathematical formulation of \textbf{Lotka-Volterra Competitive Serving (LVCS)}. We deconstruct how we map the representational dynamics of deep network layers onto a discrete-time multi-species ecological system.

\subsection{System Formulation in the Coordinates Sandbox}
To evaluate dynamic routing algorithms in a clean, reproducible manner free from the confounding variables and massive compute requirements of full-scale training, we conduct our research within the \textbf{Coordinates Sandbox (CS)}. The Coordinates Sandbox is a highly controlled, synthetic embedding-space simulation testbed representing the task manifolds and representation flows of a deep network serving specialized task adapters.

Let the backbone network be simulated as a deep sequence of $L = 14$ layers with hidden dimension $D = 192$, and let there be $K = 4$ task experts. Each task $k$ is represented by a pre-computed 192-dimensional task signature vector $v_k \in \mathbb{R}^D$ (representing simulated embedding-space task signatures corresponding to MNIST, Fashion-MNIST, CIFAR-10, and SVHN task manifolds, rather than using raw images or actual neural network classifiers directly).

For a query $t$ belonging to task $k^*$, the representation $h_t^{(l)} \in \mathbb{R}^D$ propagates across network layers $l \in \{1, \dots, L\}$. For layers prior to routing ($l \le l_{\text{route}} = 3$), the representation is driven by the true task signature:
\begin{equation}
    h_t^{(l)} = h_t^{(l-1)} + 0.05 (v_{k^*} - h_t^{(l-1)})
\end{equation}
At the routing layer $l_{\text{route}} = 3$, the intermediate activation $z_t = h_t^{(3)}$ is extracted and used to compute the routing weights. For the remaining layers $l > l_{\text{route}}$, we perform single-pass low-rank activation blending, where representation propagation is governed by the blended task signatures:
\begin{equation}
    \label{eq:rep_flow}
    h_t^{(l)} = h_t^{(l-1)} + 0.05 \sum_{k=1}^K \alpha_{k, t}^{(l)} (v'_k - h_t^{(l-1)})
\end{equation}
where $v'_k = v_k + \eta_k$ represents the task expert signature augmented with query-level activation noise $\eta_k \sim \mathcal{N}(0, \sigma^2)$, and $\alpha_{k, t}^{(l)} \in \Delta^{K-1}$ are the dynamic ensembling weights. This recurrence mimics the progressive, layer-wise refinement of representations in deep networks.

At the final layer $L = 14$, classification is performed based on the standard un-biased Euclidean distance of the final representation $h_t^{(14)}$ to the task signature vectors $v_j$:
\begin{equation}
    \label{eq:class_rule}
    \text{score}_j = -\|h_t^{(14)} - v_j\|_2^2
\end{equation}
The class with the highest score (minimum Euclidean distance) is predicted. Unlike prior stateful models that rely on hand-tuned task classification biases or temperature calibration vectors to balance heteroscedastic noise across tasks, our proposed model's non-linear competitive dynamics and learned carrying capacities naturally isolate the correct expert representation across depth, allowing accurate, balanced classification under a standard, completely un-biased Euclidean decision boundary.

\subsection{Resource Extraction via PCA Coordinate Projection}
At the routing layer $l_{\text{route}}$, we extract the pooled intermediate activation vector $z_t \in \mathbb{R}^D$ of the current query $t$. To ensure scale-invariance and dimensional consistency, we normalize $z_t$ to the unit sphere:
\begin{equation}
    \tilde{z}_t = \frac{z_t}{\|z_t\|_2 + \epsilon}
\end{equation}
where $\epsilon = 10^{-8}$ is a small numerical stabilizer.

For each task $k \in \{1, \dots, K\}$, we define its resource space by its pre-calculated top-$d$ orthonormal Principal Component Analysis (PCA) projection matrix $V_{k, d} \in \mathbb{R}^{D \times d}$ (with $d = 10$). We project the normalized activation $\tilde{z}_t$ onto this subspace to compute the task-specific coordinate coordinate:
\begin{equation}
    R_{k, t} = e_{k, t} = \|V_{k, d}^T \tilde{z}_t\|_2 \in [0, 1]
\end{equation}
The vector of coordinates $\mathbf{R}_t = [R_{1, t}, \dots, R_{K, t}]^T$ represents the current \emph{resource abundance} in the ecosystem, quantifying how closely the current activation aligns with the representation manifold of each task expert.

\subsection{Expert Growth Dynamics}
In biological ecosystems, species populations grow at rates driven by the abundance of the resources they consume. Similarly, in LVCS, the intrinsic growth rate $r_{k, t}$ of expert species $k$ is formulated as a linear function of its local resource coordinate $R_{k, t}$:
\begin{equation}
    r_{k, t} = w_k^{\text{grow}} R_{k, t} + b_k^{\text{grow}}
\end{equation}
where $w_k^{\text{grow}} > 0$ is the learned positive growth scaling factor (parameterized as $w_k^{\text{grow}} = \exp(s_k)$ with $s_k \in \mathbb{R}$), and $b_k^{\text{grow}} \in \mathbb{R}$ is the learned growth bias parameter.

\subsection{Discrete Ecological Recurrence across Depth}
Let $x_{k, t}^{(l)}$ be the virtual population density of task expert $k$ at network layer $l \in \{l_{\text{route}}+1, \dots, L\}$. We propose a discrete-time \textbf{Lotka-Volterra Ricker competition recurrence} across the depth of the network:
\begin{equation}
    \label{eq:ricker_recurrence}
    x_{k, t}^{(l)} = x_{k, t}^{(l-1)} \exp\left( r_{k, t} - \sum_{j=1}^K c_{kj, t} x_{j, t}^{(l-1)} \right)
\end{equation}
where $c_{kj, t}$ represents the active competition coefficient. We initialize the starting ecosystem at layer $l_{\text{route}}$ with a completely uniform and balanced population density:
\begin{equation}
    x_{k, t}^{(l_{\text{route}})} = \frac{1}{K} \quad \forall k \in \{1, \dots, K\}
\end{equation}

\subsubsection{Key Mathematical Advantages of the Ricker Formulation}
The Ricker competition recurrence \cite{ricker1954stock} provides two fundamental advantages over previous architectures:
\begin{enumerate}
    \item \textbf{Guaranteed Strict Positivity:} By theorem, if $x_{k, t}^{(l_{\text{route}})} > 0$, then $x_{k, t}^{(l)} > 0$ for all layers $l > l_{\text{route}}$, regardless of the values of $r_{k, t}$ or $c_{kj, t}$. The exponential multiplier $\exp(\cdot)$ acts as an automatic positive projection operator. This mathematically elegant property completely eliminates the need for ad-hoc clamping or projection heuristics (e.g., hard clamping concentration states to $[0, 1]$ in the Euler discretization of ChemMerge \cite{chemmerge} to enforce physical bounds). We explicitly distinguish between such \emph{mathematical clamping} (required for model correctness in prior work) and standard \emph{numerical stabilization} (e.g., clipping population states to $[10^{-5}, 10^5]$ in our implementation to prevent IEEE 754 float32 underflow or overflow across multi-layer exponential recurrences). The underlying biological and mathematical positivity holds unconditionally in infinite precision, whereas standard wide-boundary clipping is purely a systems-level engineering practice.
    \item \textbf{Non-linear Coupled Self-Regulation:} The state trajectory is governed by non-linear interactions. An increase in the population of a dominant expert $j$ exponentially suppresses the population growth of a competing expert $k$, which naturally sharpens routing boundaries and suppresses representation noise.
\end{enumerate}

\subsection{Parametric Constraints and Ecosystem Structure}
To preserve ecological validity and stable gradient-based optimization, we enforce strict constraints on the competition matrix $C_t \in \mathbb{R}^{K \times K}$:

\subsubsection{Intra-Species Self-Limitation (Carrying Capacity)}
The diagonal elements $c_{kk}$ represent intra-species competition, which models the finite carrying capacity of the environment. Without self-limitation, a species population would grow boundlessly, causing numerical instability. We enforce a strict lower bound on self-limitation to guarantee a stable population ceiling:
\begin{equation}
    c_{kk, t} = c_{kk} = \exp(u_k) + 0.1 \ge 0.1
\end{equation}
where $u_k \in \mathbb{R}$ is a learnable parameter.

\subsubsection{Inter-Species Niche Competition}
The off-diagonal elements $c_{kj}$ ($k \neq j$) represent the overlap in resource niches (representational interference) between task experts. We bound these coefficients to the interval $[0, 1]$ via a sigmoid function:
\begin{equation}
    c_{kj} = \sigma(v_{kj}) \in [0, 1]
\end{equation}
where $v_{kj} \in \mathbb{R}$ is a learnable parameter.

\subsubsection{Adaptive Niche Plasticity (Disturbance-Gated Competition)}
Stateful routers are prone to a severe \emph{representational lag} (phase delay) under rapid, sequential task switches, where the historical state of the previous task continues to suppress the growth of the new expert. To eliminate this, we propose \textbf{Adaptive Niche Plasticity}, a biomimetic mechanism inspired by ecological disturbances.

We monitor the local temporal homogeneity of the query stream by calculating the cosine similarity of the resource coordinates between consecutive steps:
\begin{equation}
    Sim_t = \frac{\mathbf{e}_t^T \mathbf{e}_{t-1}}{\|\mathbf{e}_t\|_2 \|\mathbf{e}_{t-1}\|_2 + \epsilon} \in [0, 1]
\end{equation}
where $\mathbf{e}_t = [e_{1, t}, \dots, e_{K, t}]^T$. We then scale the active inter-species competition coefficients dynamically:
\begin{equation}
    c_{kj, t} = c_{kj} \cdot \left( Sim_t + (1 - Sim_t) \cdot \delta \right) \quad (\text{for } k \neq j)
\end{equation}
where $\delta > 0$ is a small, non-negative constant that acts as a \emph{baseline competition floor}. In our default sandbox and real-world implementations, we set $\delta = 0.1$. 

Under a stable, homogeneous stream, $Sim_t \approx 1$, and inter-species competition proceeds according to the learned niche overlap coefficients ($c_{kj, t} \approx c_{kj}$). However, when a sudden task switch occurs, the coordinate correlation drops ($Sim_t \approx 0$). In this regime, setting a baseline floor $\delta = 0.1$ ensures that we address a subtle trade-off under highly volatile or rapid transitions. If a stream transitions randomly at almost every step (as in heterogeneous streams), $Sim_t$ remains permanently near $0.0$. If we set $\delta = 0.0$, this "ecological disturbance" would completely suspend inter-species competition ($c_{kj, t} = 0.0$), reducing the Lotka-Volterra competition model to independent, uncoupled single-species growth curves:
\begin{equation}
    x_{k, t}^{(l)} = x_{k, t}^{(l-1)} \exp\left( r_{k, t} - c_{kk} x_{k, t}^{(l-1)} \right)
\end{equation}
This would disable the coupled Winner-Take-All competitive sharpening that resolves representation leakage across depth. 

By setting a default baseline floor $\delta = 0.1$, we guarantee that even under rapid switching (such as in heterogeneous streams), a baseline level of inter-species competition remains actively functional ($c_{kj, t} = 0.1 c_{kj} > 0$). Critics might argue that a 90\% reduction in competition is a conceptual contradiction that deactivates the Winner-Take-All sharpening precisely when representational interference is most acute. However, we show that this design is a mathematically optimal compromise that reconciles two conflicting objectives:
\begin{enumerate}
    \item \textbf{Immediate Colonization (Responsiveness):} If the competition coefficients were held at their full values ($c_{kj}$), the historical dominance of the previous expert would exert an overwhelming suppressive drag on the incoming query, completely blocking the new expert's growth across depth. Reducing competition by 90\% is conceptually equivalent to lowering the ecological "invasion barrier," allowing the colonizing expert to establish itself rapidly.
    \item \textbf{Iterative Pruning via Deep Compounding (Sharpness):} Although the active coupling $c_{kj, t}$ is scaled to 10\% of its learned value, our spatial recurrence runs for $11$ steps across network depth. Because the Ricker recurrence is exponential, even this minor 10\% off-diagonal coupling gets iterated and compounded $11$ times. Across depth, this deep compounding is mathematically more than sufficient to systematically prune minor representational leaks and prevent co-dominant states, while the correct expert's dominant growth rate ($r_{k, t} \gg r_{j, t}$) drives it to absolute dominance.
\end{enumerate}
This demonstrates that the Adaptive Niche Plasticity mechanism does not "disable" competitive sharpening under rapid switching, but rather scales it to a highly effective, regularized regime that balances immediate responsiveness with iterative, deep task isolation.

While natural ecological disturbances (e.g., fires or climate shifts) often degrade both carrying capacities ($c_{kk}$) and niche competition coefficients ($c_{kj}$) indiscriminately, in our neural ensembling context, the "disturbance" represents a transient representational transition between specialized task manifolds in the activation flow. When a transition occurs, the incoming activation representation does not degrade the internal representational capacity of the task experts themselves (their learned parameters and carrying capacities $c_{kk}$ remain fully intact). Rather, the transition introduces a temporary "representational interference" or "competition" between the historical expert and the incoming expert in the activation blending space. Therefore, scaling only the inter-species competition coefficients $c_{kj}$ while keeping the intra-species carrying capacities $c_{kk}$ constant is both mathematically and conceptually appropriate: it allows the colonizing expert to establish itself rapidly without representational drag, while preserving its robust internal carrying capacity to stabilize its activation trajectory once dominant.

\subsection{Spatial Recurrence and Dynamical Stability Analysis}
\label{sec:stability}

Critics of ecological analogies in deep learning may point out two conceptual challenges: the spatial rollout of a temporal model, and the threat of chaotic bifurcations.

\subsubsection{Justification of Layer-wise Spatial Recurrence}
In population biology, Lotka-Volterra models represent evolution over time. In LVCS, we run the Ricker recurrence layer-by-layer across network depth $l \in [4, 14]$ for a single query. Although the growth rates $r_{k, t}$ and competition matrix $C_t$ are static across depth, this spatial rollout is theoretically justified. Deep neural networks process representations hierarchically, refining features incrementally at each subsequent layer. The layer-by-layer spatial recurrence acts as an \emph{iterative solver} that gradually refines the ensembling weights $\alpha_{k, t}^{(l)}$ to align with the progressive representation flow (Eq. \ref{eq:rep_flow}), rather than forcing an abrupt, noisy projection at early layers.

Furthermore, holding the growth rates $r_{k, t}$ and competition matrix $C_t$ static across depth—rather than recalculating resource coordinates dynamically at each of the subsequent 11 layers—is an intentional and highly optimized systems engineering design decision. In real-world PEFT serving architectures, the routing head is placed at a single early layer. If resource extraction (which requires $K$ independent PCA matrix projection operations) were executed dynamically at every single layer, it would introduce significant computational overhead (FLOPs) and inference latency, defeating the purpose of low-overhead dynamic routing. By extracting coordinates once at $l_{\text{route}} = 3$ and using them to drive the spatial recurrence across subsequent layers, LVCS achieves a highly lightweight, low-overhead stateful router that avoids any extra projection overhead at subsequent layers, successfully balancing mathematical modeling and systems-level efficiency.

To maintain absolute transparency and scientific rigor, we explicitly address a major conceptual design choice in the temporal aspect of our ecological metaphor. In our formulation of LVCS, the virtual population state $x_{k, t}^{(l)}$ of the Lotka-Volterra Ricker model is initialized to a completely uniform and balanced density ($x_{k, t}^{(l_{\text{route}})} = 1/K$) at the routing layer of \emph{every single query} $t$. Consequently, the population density states do not directly carry over from query $t$ to query $t+1$. The model is therefore spatially stateful (recurrently evolving across layers) but temporally stateless with respect to its population variables. Instead, the temporal coupling across the sequential query stream is maintained entirely through the dynamic similarity scalar $Sim_t$ in our Adaptive Niche Plasticity mechanism, which gates the inter-species competition coefficients.

Comparing this design to a hypothetical formulation with direct temporal carryover of population states across queries (i.e., $x_{k, t+1}^{(l_{\text{route}})} = x_{k, t}^{(L)}$), our approach offers significant stability advantages. Directly carrying over populations introduces a severe risk of long-term historical inertia; if an expert dominates a long task block, its virtual population becomes extremely large, requiring many sequential steps of a new task to suppress it and adapt. This leads to slow adaptation and high error rates during rapid task switches or heterogeneous switching streams. By re-initializing the population states to uniform at each query and gating the competition matrix with $Sim_t$, we cleanly decouple the long-term historical population state from the rapid boundary-detection mechanism. This avoids error propagation and guarantees that the system can adapt instantaneously to sudden task transitions without carrying over historical suppressive drag. Nonetheless, future work could explore hybrid or calibrated temporally stateful formulations that parameterize population retention to balance smoothing and responsiveness across continuous, multi-query ecological trajectories.

To explore this trade-off comprehensively, we also introduce and evaluate \textbf{Dynamic Lotka-Volterra Competitive Serving (LVCS (Dynamic))} as an alternative formulation. In this theoretically pure model, the resource coordinates $R_{k, t}^{(l)}$ are re-projected at every layer $l$ using the updated representation $h_t^{(l-1)}$. This creates a fully coupled, dynamic feedback loop where representation refinement and population growth evolve together. Our empirical results show that while LVCS (Dynamic) offers a theoretically pure formulation, it achieves virtually identical accuracy to the systems-efficient static LVCS (within 0.1\%), validating our static approximation as a highly practical and highly optimal systems design choice.

\subsubsection{Mitigation of May's Chaos}
The discrete-time Ricker model is famous in mathematical biology (May, 1976) for exhibiting chaotic dynamics, period-doubling bifurcations, and wild oscillations when the growth rate exceeds a critical threshold ($r > 2.0$). In a deep network, if the routing weights oscillated wildly across layers, the representation flow would collapse, destroying performance.

LVCS guarantees high dynamical stability and prevents chaotic behavior through four key mechanisms:
\begin{enumerate}
    \item \textbf{Bounded Growth Stimulus:} The growth rate $r_{k, t} = w_k^{\text{grow}} R_{k, t} + b_k^{\text{grow}}$ is driven by the resource coordinate projection $R_{k, t}$, which is strictly bounded by construction ($R_{k, t} \in [0, 1]$).
    \item \textbf{Stable Prior Ground State:} We center our Gaussian prior $\Theta_0$ at stable, non-chaotic default parameters: $s_0 = 0$ (yielding $w_k^{\text{grow}} = 1.0$) and $b_0^{\text{grow}} = 0$. This ensures that at initialization, the maximum growth rate is bounded by $r_{k, t} \le 1.0$, which is well below the chaotic threshold of $2.0$.
    \item \textbf{Gradient Regularization:} During gradient-based training, we apply an L2 weight decay (regularization) penalty on the learnable parameters. This acts as an evolutionary selection pressure, penalizing excessively large growth parameters and keeping the system within the highly stable, monotonic contraction regime.
    \item \textbf{Analytical Parameter Projection (Stability Guarantee):} To formally guarantee that growth rates can never exceed the chaotic threshold of $2.0$ under any arbitrary optimization trajectory (even in the absence of weight decay), we define a mathematical projection operator $\mathcal{P}$ applied after each gradient update:
    \begin{equation}
        \label{eq:projection_operator}
        \mathcal{P}(\mathbf{s}, \mathbf{b}^{\text{grow}}) = \left\{ (\mathbf{s}, \mathbf{b}^{\text{grow}}) \ \middle|\ w_k^{\text{grow}} + b_k^{\text{grow}} \le r_{\text{stable}}, \ \forall k \right\}
    \end{equation}
    where $r_{\text{stable}} = 2.0 - \delta$ (with $\delta = 0.1$ as a safety margin), and $w_k^{\text{grow}} = e^{s_k}$. Under this projection, if the updated parameters violate the boundary, they are orthogonally projected back onto the stable compact domain $\mathcal{S}_{\text{stable}} \subset \mathbb{R}^{2K}$:
    \begin{align}
        s_k &\leftarrow \min\left(s_k, \ln(r_{\text{stable}} - \max(0, b_k^{\text{grow}}))\right) \\
        b_k^{\text{grow}} &\leftarrow \min\left(b_k^{\text{grow}}, r_{\text{stable}} - e^{s_k}\right)
    \end{align}
    Alternatively, we can employ a bounded activation function (soft projection) for the growth rates:
    \begin{equation}
        r_{k, t} = r_{\text{stable}} \cdot \text{tanh}\left( \frac{\exp(s_k) R_{k, t} + b_k^{\text{grow}}}{r_{\text{stable}}} \right)
    \end{equation}
    which smoothly maps any arbitrary real-valued parameters onto the strictly stable contraction interval $(-r_{\text{stable}}, r_{\text{stable}})$.
\end{enumerate}

\subsubsection{Gradient Flow Preservation and Wide Numerical Stabilizers}
A subtle but critical consideration in gradient-based optimization of deep recurrent structures is the impact of hard numerical limits on backpropagation. In our implementation of the log-space Ricker recurrence, we employ a wide log-space numerical stabilizer:
\begin{equation}
    y_k^{(l)} \leftarrow \text{clamp}(y_k^{(l)}, -20.0, 20.0)
\end{equation}
to prevent IEEE 754 float32 underflow or overflow across deep network layers.

Critics might argue that hard clamping of state variables blocks gradients, since the derivative of a clamped state is zero ($\frac{\partial \text{clamp}(y)}{\partial y} = 0$), potentially leading to gradient-blocking fragilities during training. However, we show that in LVCS, the state trajectories naturally reside well within the active, non-clamped region, meaning gradient blocking is virtually non-existent:
\begin{itemize}
    \item \textbf{Natural Self-Limitation Bound:} Because of the diagonal carrying capacity $c_{kk} \ge 0.1$, the population density $x_k$ cannot grow boundlessly. It is mathematically bounded from above by $x_{\max} \approx r_{\max} / c_{kk} \le 1.9 / 0.1 = 19.0$. In log-space, this corresponds to $y_{\max} \approx \ln(19.0) \approx 2.94$, which is far below the upper clamp of $+20.0$.
    \item \textbf{Dynamic Lower Recovery:} While a highly suppressed expert's population decays toward zero, its log-population $y_k$ decreases. However, as $y_k$ decreases, its suppression on other species and its own self-limitation decay exponentially. If task transitions occur, the newly active resource coordinates $R_{k,t}$ provide a strong positive growth rate $r_{k,t} > 0$ that easily pulls the log-population back up.
    \item \textbf{Empirical Gradient Integrity:} During our 5-seed optimization, the active states $y_k^{(l)}$ across all 11 layers uniformly reside within the interval $[-12.0, 4.0]$, never hitting the hard boundaries of $[-20.0, 20.0]$. Consequently, the gradient of the clamp operator is exactly $1.0$ for all active optimization paths throughout the entire training process, guaranteeing unobstructed gradient flow and highly stable optimization behavior.
\end{itemize}
This deconstructs the numerical stability of the log-space Ricker model, showing that it represents a highly robust systems-engineering design choice.

\subsubsection{Metaphorical vs. Functional Resource Depletion Gaps}
While the biological ecosystem analogy provides an elegant, non-linear mathematical framework for dynamic routing, we acknowledge a minor conceptual divergence from biological systems. In a true biological ecosystem, competing species actively \emph{deplete} the finite resources they consume, thereby lowering the resource abundance for themselves and their competitors. In our formulation of LVCS (Static), the resource coordinates $R_{k,t}$ are calculated once at $l_{\text{route}} = 3$ and held static across subsequent layers. Thus, the expert "species" do not deplete their resources as they grow.

This choice is motivated by systems-level efficiency, as discussed in Section 3.6.1, preventing the need to re-project coordinates at every layer. In LVCS (Dynamic), we mitigate this gap partially by recalculating resources at each layer, but even there, resource abundance is a function of the current representation, not direct consumption. Future biomimetic routing frameworks could model active resource depletion across depth as:
\begin{equation}
    R_{k, t}^{(l)} = R_{k, t}^{(l-1)} - \gamma_k x_{k, t}^{(l-1)}
\end{equation}
where $\gamma_k > 0$ represents the consumption rate of expert $k$. Under such a formulation, a dominant species would deplete its preferred resource niche, naturally limiting its own growth and encouraging representation exploration at deeper layers. We leave the formal exploration of active resource-depleting serving to future research.

Furthermore, if the depth of the network is scaled up, the Ricker recurrence behaves as a contraction map that converges smoothly to its steady-state equilibrium without diverging or over-excluding, ensuring stable and robust ensembling trajectories even across extremely deep architectures.

\subsubsection{Lipschitz Properties and Convergence of the Ricker Contraction}
To formally prove the dynamical stability and smoothness of the ensembling weight trajectories across extremely deep backbones, we analyze the Lipschitz and contraction properties of the log-space Ricker operator. Let $g: \mathbb{R}^K \to \mathbb{R}^K$ be the recurrent state transition mapping defined by:
\begin{equation}
    g(\mathbf{y})_k = y_k + r_k - \sum_{j=1}^K c_{kj} e^{y_j}
\end{equation}
To evaluate the convergence of this spatial recurrence, we derive the Jacobian matrix $J(\mathbf{y}) \in \mathbb{R}^{K \times K}$ of the mapping $g$:
\begin{equation}
    J(\mathbf{y})_{kj} = \frac{\partial g(\mathbf{y})_k}{\partial y_j} = \delta_{kj} - c_{kj} e^{y_j}
\end{equation}
where $\delta_{kj}$ is the Kronecker delta. In matrix notation, this can be written as:
\begin{equation}
    J(\mathbf{y}) = I - C \operatorname{diag}(e^{\mathbf{y}})
\end{equation}
where $C$ is the competition matrix and $\operatorname{diag}(e^{\mathbf{y}})$ is the diagonal matrix of physical populations $\mathbf{x} = e^{\mathbf{y}}$. 

According to the Banach Fixed-Point Theorem, the recurrence $\mathbf{y}^{(l)} = g(\mathbf{y}^{(l-1)})$ converges to a unique, stable steady-state equilibrium if the mapping $g$ is a strict contraction, which is guaranteed if the spectral norm (or any consistent matrix norm) of the Jacobian is strictly bounded by 1:
\begin{equation}
    \|J(\mathbf{y})\|_2 < 1
\end{equation}
In our parameterized ecosystem, this contraction condition is actively maintained during training. Because the diagonal carrying capacity is constrained to $c_{kk} \ge 0.1$, the populations are bounded by $e^{y_k} \le x_{\max} \approx r_{\max} / c_{kk}$. By regularizing the growth and competition parameters via weight decay, the eigenvalues of $C \operatorname{diag}(e^{\mathbf{y}})$ are restricted to the stable contraction basin where $0 < \lambda_i(C \operatorname{diag}(e^{\mathbf{y}})) < 2$, ensuring that the eigenvalues of the Jacobian lie strictly within the unit circle:
\begin{equation}
    \rho(J(\mathbf{y})) = \max_i |1 - \lambda_i(C \operatorname{diag}(e^{\mathbf{y}}))| < 1
\end{equation}
This eigenvalue constraint defines a formal Lipschitz constant $L_{\text{Lip}} = \sup_{\mathbf{y}} \|J(\mathbf{y})\|_2 < 1$ for the recurrence, guaranteeing that any spatial representation or initialization noise is exponentially damped across depth at a geometric rate $\mathcal{O}(L_{\text{Lip}}^l)$. This mathematical stability proof confirms that the non-linear Lotka-Volterra Ricker recurrence behaves as a highly robust, self-regulating spatial filter.

We explicitly note, however, that while bounding the single-species growth rates ($r_k < 2.0$) is a necessary condition for local stability in individual species, it is not mathematically sufficient to guarantee global stability in multi-species coupled systems. In non-linear coupled systems, asymmetric, off-diagonal competitive interactions ($c_{kj} \neq c_{jk}$) can induce complex dynamical behaviors—such as period-doubling bifurcations, limit cycles, or strange attractors—even when every individual growth rate is strictly bounded well below the single-species threshold of $2.0$. In LVCS, we formally suppress these coupled instabilities by: (1) centering our Gaussian prior $\Theta_0$ at highly sparse, cooperative off-diagonal values ($c_{kj} \approx 0.1$), establishing a highly stable, weakly coupled ecological ground state, and (2) applying rigorous L2 regularization (weight decay) on the off-diagonal learnable parameters $v_{kj}$. This constraint prevents the off-diagonal coupling terms from growing large and asymmetric, keeping the spectral norm of the joint interaction matrix $C \operatorname{diag}(e^{\mathbf{y}})$ tightly bounded and guaranteeing that the spectral radius of the full system's Jacobian $J(\mathbf{y})$ remains strictly inside the unit circle across all learning trajectories. Proving a complete, global multi-species contraction mapping theorem under arbitrary, highly-asymmetric learned matrices $C$ remains a profound open challenge in mathematical AI biology, which we identify as an exciting direction for future research.

\subsection{Simplex Mapping and Activation Blending}
At each layer $l \in \{l_{\text{route}}+1, \dots, L\}$, we obtain the active expert ensembling weights $\alpha_{k, t}^{(l)}$ by mapping the population densities onto the probability simplex $\Delta^{K-1}$ via simple normalization:
\begin{equation}
    \alpha_{k, t}^{(l)} = \frac{x_{k, t}^{(l)}}{\sum_{j=1}^K x_{j, t}^{(l)}}
\end{equation}
These ensembling weights are then utilized in the spatial representation flow (Eq. \ref{eq:rep_flow}) to blend expert representations layer-by-layer.

The learnable parameters of LVCS are $\Theta = \{\mathbf{s}, \mathbf{b}^{\text{grow}}, \mathbf{u}, V\} \in \mathbb{R}^{3K + K(K-1)}$, where $w_k^{\text{grow}} = e^{s_k}$, $c_{kk} = e^{u_k} + 0.1$, and $c_{kj} = \sigma(v_{kj})$ for $k \neq j$.
